\begin{afterword}
\summaryhead

The post starts from Robin Hanson's proposal to allow stores that sell banned products. The author had replied that some poor, poorly educated mother of five would buy a poison there and die, and some readers took this as an argument for regulation. The author calls it a factual observation.

On a question of fact the evidence should end up on one side. A policy with many consequences has no reason to be one-sided, yet people want their policy debates to be, because arguments are soldiers and every cost of one's own side is an enemy. The reverse mistake, splitting the difference, is also rejected.

People also dislike an unfair universe, and deny either the facts or the unfairness: some say the buyer deserved her death. The post answers with the birth lottery of genes and upbringing. Real tough-mindedness admits that the mother did not deserve to die and keeps the shops open anyway if the cost-benefit calculation says so. The author counts such deaths as tragedies.

\responsehead

The thesis is right, and the post gives a reason for it. A question of fact has one answer, so strong evidence ends up favoring one side. A policy has many consequences, so there is no reason for all its costs to fall on one side. The post also rejects the reverse mistake: some policies really are lopsided, and splitting the difference is not wisdom. It applies its charge to libertarians and to believers in regulation alike.

What the post does not support is its account of why people want one-sided debates. The account is the image of arguments as soldiers, repeated almost word for word from ``Politics is the Mind-Killer,'' where it was also given without evidence. The tendency to ``deny all costs of a favored policy'' is asserted. The closest evidence, the study the author cites ten days later, shows that people judge risks lower when they judge benefits higher. That is a bias in judgment, not a denial of every cost.

The example has two problems. The post asks why readers took the dying mother as an argument for regulation, but does not mention that the author's reply was worded as a rebuttal: it began ``But even so,'' after three arguments for the shops. And the claim that if the buyer deserved her fate ``there would be no downside'' to the shops does not follow from the post's own story, whose first paragraph included her orphans.

Two factual points are inaccurate. The heritability of intelligence is given as 0.6 to 0.8 on ``overwhelming'' evidence, where summaries give 0.4 to 0.8, lower in children; the post's argument does not depend on the number. The footnote for Hanson's proposal cites a debate published three months after the post, which does not contain it. The proposal itself, posted the day before, required buyers to pass a test showing that they knew regulators disapproved, a detail the post leaves out.

In short: a sound point about trade-offs, argued well, with an unsupported account of why people resist it, and an example whose puzzle, why readers heard an argument for regulation, is explained by the wording of the author's own reply.
\end{afterword}
